% !TEX root = main.tex
We express the ring-switching construction of Diamond and
Posen~\cite{eurocrypt/DiamondP26} in our vector and matrix notation
to support the weighted tensor queries of Section~\ref{sec:blocks}.

Let \(N=m\kappa\) and \([\ExtensionField:\Field]=\kappa\).
We commit to the fixed-basis packing \(\bm A\ext{\bm\beta}\)
of \(\bm A\in\Field^{m\times\kappa}\) and reduce each query
to one MLE opening. The verifier evaluates the resulting public
query by a recursion on small base-field matrices.
We then instantiate the underlying PCS from a base-field linear
code via Blaze's code-switching construction.

\subsubsection{Given a PCS over the Extension Field}

\medskip\noindent\textbf{Commitment and opening relation.}

Let \(m,\kappa\) be powers of two, with \(\kappa\geq2\), and write
\(\ell=\log_2m\).
Fix an \(\Field\)-basis
\(\ext{\bm\beta}=(\ext\beta_1,\ldots,\ext\beta_\kappa)^\top\)
of \(\ExtensionField\), with \(\ext\beta_1=1\), and let \(\pc\)
be a multilinear PCS for vectors in \(\ExtensionField^m\).
Extension-field elements are stored in this basis, so packing
and coordinate extraction only reorganize their base-field coordinates.
The schoolbook cost bounds below use a polynomial basis.

\paragraph{Setup and commitment.}
Run \(\pc.\Setup(1^\lambda,m)\) to obtain
\(\mathsf{pp}_{\pc}\), and include the fields, basis, and dimensions
\(N,m,\kappa\) in the public parameters.
For the witness matrix
\(\bm A=\concat{\bm a^{(1)},\ldots,\bm a^{(\kappa)}}
\in\Field^{m\times\kappa}\), pack each row into one
extension-field element, obtaining
\(\ext{\bm a}=\bm A\ext{\bm\beta}\in\ExtensionField^m\).
Compute the commitment
\(\mathsf{cm}=\pc.\Com(\mathsf{pp}_{\pc},\ext{\bm a};
\rho_{\pc})\), where \(\rho_{\pc}\) is the commitment
randomness, and retain \(\bm A\) and \(\rho_{\pc}\).
The fixed basis makes binding to \(\ext{\bm a}\) equivalent
to binding to \(\bm A\).

\paragraph{Opening relation.}
A query specifies column weights
\(\ext{\bm\rho}\in\ExtensionField^\kappa\), a tensor vector
\(\ext{\bm\xi}\in\ExtensionField^m\), and a claimed value
\(\ext v\in\ExtensionField\). The opening relation is:
\begin{equation}
    \ext v
    =\langle\bm A\ext{\bm\rho},\ext{\bm\xi}\rangle
    =\ext{\bm\xi}^{\top}\bm A\ext{\bm\rho}.
\label{eq:pcs-opening-relation}
\end{equation}
The row query is supplied as
\(\ext{\bm\xi}=\bigotimes_{j=1}^{\ell}\ext{\bm\xi}^{(j)}\),
where each factor is
\(\ext{\bm\xi}^{(j)}=(\ext\xi^{(j)}_0,
\ext\xi^{(j)}_1)^\top\in\ExtensionField^2\).
Ordinary MLE queries use the tensor spans of the corresponding
parts of the evaluation point for both \(\ext{\bm\rho}\) and
\(\ext{\bm\xi}\). Allowing arbitrary two-entry factors also
supports the coordinatewise products of tensor queries arising
in the weighted Hadamard reduction.

\medskip\noindent\textbf{Packed opening protocol.}

The commitment uses the fixed weights \(\ext{\bm\beta}\), whereas
the query uses \(\ext{\bm\rho}\). The prover first sends the
column evaluations. Coordinate decomposition and a random
projection then reduce their correctness to an inner-product
claim about the committed vector.

\paragraph{Reduce the column claims.}
Write the public query and the claimed column values in base-field
coordinates:
\[
    \ext{\bm\xi}=\bm\Xi\ext{\bm\beta},
    \qquad
    \ext{\bm u}^{\top}=\ext{\bm\beta}^{\top}\bm U,
    \qquad
    \bm\Xi\in\Field^{m\times\kappa},
    \quad
    \bm U\in\Field^{\kappa\times\kappa}.
\]
The honest prover computes \(\bm U=\bm\Xi^\top\bm A\) and sends
\(\ext{\bm u}=\bm U^\top\ext{\bm\beta}\).
The verifier extracts \(\bm U\) from this message and checks
\(\ext v=\ext{\bm u}^{\top}\ext{\bm\rho}\).
The matrix \(\bm\Xi\) is defined implicitly by the public query.

Since \(\bm A\) has base-field entries, the remaining claim
\(\ext{\bm u}^{\top}=\ext{\bm\xi}^{\top}\bm A\) is equivalent
to \(\bm U=\bm\Xi^\top\bm A\). Packing the rows of this
identity introduces the committed vector \(\ext{\bm a}\):
\begin{equation}
    \bm U=\bm\Xi^\top\bm A
    \quad\Longleftrightarrow\quad
    \bm U\ext{\bm\beta}
    =\bm\Xi^\top\ext{\bm a}.
\label{eq:pcs-encoded-matrix}
\end{equation}
After receiving \(\ext{\bm u}\), the verifier samples
\(\ext{\bm r}\sample\ExtensionField^\kappa\), independently of
the original query weights, and projects
Eq.~\eqref{eq:pcs-encoded-matrix} to one inner-product claim:
\begin{equation}
    \ext c
    =\ext{\bm r}^{\top}\bm U\ext{\bm\beta}
    =\langle\ext{\bm a},\bm\Xi\ext{\bm r}\rangle.
\label{eq:pcs-random-opening}
\end{equation}
The verifier computes \(\ext c\) by packing the rows of
\(\bm U\) and taking their inner product with \(\ext{\bm r}\).

Apply Fact~\ref{fact:ip} over \(\ExtensionField\) to
Eq.~\eqref{eq:pcs-random-opening}.
Let \(\ext{\bm y}\in\ExtensionField^\ell\) be its terminal point.
The prover's terminal values
\(\ext w_{\mathsf{com}},\ext w\in\ExtensionField\) give two claims:
\begin{equation}
    \ext w_{\mathsf{com}}
    =\langle\ext{\bm a},\tensorspan{\ext{\bm y}}\rangle,
    \qquad
    \ext w=\tensorspan{\ext{\bm y}}^{\top}\bm\Xi\ext{\bm r}.
\label{eq:pcs-terminal-claims}
\end{equation}
The first is an MLE opening of the committed vector.
The verifier checks the second by the following matrix recursion.

\paragraph{Evaluate the public query.}
Write \(\tensorspan{\ext{\bm y}}=\bm Y\ext{\bm\beta}\), where
\(\bm Y\in\Field^{m\times\kappa}\).
It suffices to compute the small matrix:
\begin{equation}
    \bm W=\bm Y^\top\bm\Xi\in\Field^{\kappa\times\kappa},
    \qquad
    \ext w=\ext{\bm\beta}^{\top}\bm W\ext{\bm r}.
\label{eq:pcs-query-matrix}
\end{equation}
The public tensor factors determine \(\bm Y\) and \(\bm\Xi\),
so the verifier can compute their product without forming either
length-\(m\) coordinate table.

For \(\ext z\in\ExtensionField\), let
\(\bm M_{\ext z}\in\Field^{\kappa\times\kappa}\) be the coordinate
matrix of multiplication by \(\ext z\), characterized by:
\begin{equation}
    \ext{\bm\beta}^{\top}\bm M_{\ext z}
    =\ext z\,\ext{\bm\beta}^{\top}.
\label{eq:pcs-multiplication-matrix}
\end{equation}
Thus left multiplication by \(\bm M_{\ext z}\) scales each
coordinate column by \(\ext z\) in \(\ExtensionField\), and right
multiplication by \(\bm M_{\ext z}^\top\) scales each coordinate row.
In particular, \(\bm M_{1-\ext z}=\bm I-\bm M_{\ext z}\).

Let \(\bm Y_j,\bm\Xi_j\) be the coordinate matrices for the
suffix factors \(j,\ldots,\ell\):
\[
    \bm Y_j\ext{\bm\beta}
    =\bigotimes_{i=j}^{\ell}
        \begin{pmatrix}1-\ext y_i\\ \ext y_i\end{pmatrix},
    \qquad
    \bm\Xi_j\ext{\bm\beta}
    =\bigotimes_{i=j}^{\ell}\ext{\bm\xi}^{(i)}.
\]
Here \(\bm Y_1=\bm Y\), \(\bm\Xi_1=\bm\Xi\), and
\(\bm Y_{\ell+1}=\bm\Xi_{\ell+1}=\bm e_1^\top\), where
\(\bm e_1=(1,0,\ldots,0)^\top\in\Field^\kappa\).
The two possible values of the first suffix bit give the block
decompositions:
\begin{equation}
\begin{aligned}
    \bm Y_j
    &=
    \begin{bmatrix}
        \bm Y_{j+1}\bm M_{1-\ext y_j}^\top\\
        \bm Y_{j+1}\bm M_{\ext y_j}^\top
    \end{bmatrix},\\
    \bm\Xi_j
    &=
    \begin{bmatrix}
        \bm\Xi_{j+1}\bm M_{\ext\xi^{(j)}_0}^\top\\
        \bm\Xi_{j+1}\bm M_{\ext\xi^{(j)}_1}^\top
    \end{bmatrix}.
\end{aligned}
\label{eq:pcs-query-blocks}
\end{equation}
Consequently, the verifier computes
\(\bm W_j=\bm Y_j^\top\bm\Xi_j\) by initializing
\(\bm W_{\ell+1}=\bm e_1\bm e_1^\top\) and processing
\(j=\ell,\ldots,1\):
\begin{equation}
\begin{aligned}
    \bm W_j
    &=
    \bm M_{1-\ext y_j}\bm W_{j+1}
        \bm M_{\ext\xi^{(j)}_0}^\top
    +\bm M_{\ext y_j}\bm W_{j+1}
        \bm M_{\ext\xi^{(j)}_1}^\top\\
    &=
    \bm W_{j+1}\bm M_{\ext\xi^{(j)}_0}^\top
    +\bm M_{\ext y_j}\bm W_{j+1}
        \bm M_{\ext\xi^{(j)}_1-\ext\xi^{(j)}_0}^\top.
\end{aligned}
\label{eq:pcs-query-recurrence}
\end{equation}
This recursion yields \(\bm W=\bm W_1\) without forming
the suffix coordinate matrices.
Each step of the second line uses at most \(3\kappa\)
extension-field multiplications, implemented by row and column
scaling without materializing \(\bm M_{\ext z}\).
The recursion takes \(O(\kappa\ell)\) extension-field
multiplications and \(O(\kappa^2\ell)\) base-field additions and
coordinate moves, using \(O(\kappa^2)\) base-field coordinates
of working memory. It applies in every characteristic.

The verifier packs the columns of \(\bm W\), computes
\(\ext{\bm\beta}^{\top}\bm W\ext{\bm r}\), and checks the second
identity in Eq.~\eqref{eq:pcs-query-matrix}.
This computation is exact and requires no additional prover message.

\paragraph{Terminal opening.}
Finally, the parties authenticate \(\ext w_{\mathsf{com}}\)
at \(\ext{\bm y}\) by invoking
\(\pc.\Open(\mathsf{pp}_{\pc},\mathsf{cm},
\ext w_{\mathsf{com}},\ext{\bm y}
:\ext{\bm a},\rho_{\pc})\).
The verifier accepts if the scalar checks, the inner-product
reduction, and the underlying opening accept.
Algorithm~\ref{alg:packed-pcs} in Appendix~\ref{appx:pseudocode}
summarizes the protocol.

\medskip\noindent\textbf{Security and complexity.}

\begin{proposition}[Packed polynomial commitment scheme]
\label{proposition:packed-pcs}
Let \(\pc\) bind the prover to a vector in \(\ExtensionField^m\),
with evaluation-soundness error \(\epsilon_{\mathsf{PCS}}\) for
adaptively chosen MLE claims about that vector.
The packed commitment is binding and the opening protocol is
complete, with evaluation-soundness error at most:
\begin{equation}
    \epsilon_{\mathsf{PCS}}
    +\frac{1+2\ell}{|\ExtensionField|}.
\label{eq:pcs-soundness}
\end{equation}
The opening reduction requires no additional lower bound on
\(|\Field|\).

Excluding the underlying PCS opening, the prover's arithmetic cost is:
\begin{equation}
    O(\kappa N)\,\bb
    +O(N)\,\be
    +O(m+\kappa)\,\ee.
\label{eq:pcs-prover-cost}
\end{equation}
The verifier's arithmetic cost is:
\begin{equation}
    O\bigl(\kappa^2(1+\ell)\bigr)\,\bb
    +O\bigl(\kappa(1+\ell)\bigr)\,\ee.
\label{eq:pcs-verifier-cost}
\end{equation}
The verifier bound includes base-field additions and coordinate
reorganization in the matrix recursion.
When \(m\geq\kappa\), schoolbook arithmetic gives
\(O(\kappa N)\) base-field operations for the prover.
Peak prover memory is \(O(N+\kappa^2)\) base-field coordinates,
or \(O(N)\) when \(m\geq\kappa\), excluding the workspace of the
underlying PCS.

With \(\ext v\) already part of the public statement,
prover-to-verifier communication, excluding the underlying
opening, consists of \(\kappa+2\ell+2\) extension-field elements.
The packed commitment inherits hiding from the underlying
commitment. The opening protocol reveals the column evaluations
\(\ext{\bm u}\), so this inheritance does not establish
zero knowledge of the opening transcript.
\end{proposition}

The proof is deferred to Appendix~\ref{appx:proof-pcs}.

\subsubsection{Given a Linear Code over the Base Field}

A base-field generator matrix commutes with basis packing and
therefore also encodes packed messages over \(\ExtensionField\).
We instantiate the preceding construction via Blaze's code-switching
protocol~\cite{eurocrypt/BrehmCFRRZ25}, with a BaseFold
code~\cite{crypto/ZeilbergerCF24} as the target, under Blaze's
source-encoding and proximity assumptions.

Let \(\Code:\Field^m\to\Field^n\) be a linear code with
generator matrix \(\bm G\in\Field^{n\times m}\).
Its \(\kappa\)-way interleaving encodes the columns of
\(\bm A\in\Field^{m\times\kappa}\) independently.
Packing the resulting codeword in the fixed basis gives:
\begin{equation}
\begin{aligned}
    \bm C&=\bm G\bm A\in\Field^{n\times\kappa},\\
    \ext{\bm c}&=\bm C\ext{\bm\beta}
       =\bm G(\bm A\ext{\bm\beta})
       =\Code_{\ExtensionField}(\ext{\bm a}).
\end{aligned}
\label{eq:pcs-code-packing}
\end{equation}
Here \(\Code_{\ExtensionField}\) is the scalar extension defined
by the same generator matrix; it has the same rate and relative
distance as \(\Code\).
Thus a row of \(\bm C\) represents one extension-field codeword
symbol. Committing to these rows supplies Blaze's source codeword
oracle for \(\ext{\bm a}=\bm A\ext{\bm\beta}\).

Apply Blaze's code-switching construction over \(\ExtensionField\)
to this oracle, switching to a BaseFold encoding of the same
length-\(m\) message \(\ext{\bm a}\).
The resulting MLE openings instantiate the underlying PCS \(\pc\)
in the preceding construction.
The packed opening reduction handles the challenge-dependent weights
\(\ext{\bm\rho}\), and code switching authenticates its terminal
fixed-basis opening. We include the costs and soundness error of
the chosen Blaze instantiation, including source-encoding
consistency, in the underlying PCS costs and
\(\epsilon_{\mathsf{PCS}}\), respectively.

If \(\Code\) is constant-rate and linear-time encodable, forming
\(\bm C\) costs \(O(N)\) base-field operations and its payload
occupies \(O(N\log_2|\Field|)\) bits, excluding Merkle-tree hashes.
The BaseFold stage works on \(m=N/\kappa\) extension-field message
symbols. Thus the source encoding retains the witness's small-field
representation, and the extension-field opening operates on the
shorter packed message.
